\chapter{Safety and First Aid}
\section{Get Rolling}
\textbf{Tick the pictures in which you think the child or children are safe.}
\section{What to Know}
Fire safety, internet safety, and first aid for cuts and wounds, sprains, fractures, nosebleeds, burns and animal bites
\section{\textbf{Learning Objectives}}
\begin{itemize}
    \item to identify safe and unsafe situations
    \item to demonstrate how to protect oneself during fire
    \item to describe various fire safety rules
    \item to define first aid
    \item to apply first aid techniques for different emergency situations
\end{itemize}
Be it home, school, road or playground, we should always be careful and act responsibly. Sometimes, even a little careless attitude can lead to major accidents. So, it is very important for all of us to follow safety rules and stay safe.
\vspace{0.6in}
\section{I Explore}
\section{Fire Safety}
Fires can result from electric wire faults, gas leaks, fireworks and accidents with flammable substances.
\subsection{Safety Tips}
\begin{enumerate}
    \item Avoid playing with matches or lighters;let adults use them.
    \item Report a fire or smoke to adults immediately.
    \item Stay away from gas stoves and never play near them.
    \item Install smoke detectors and respond seriously to alarms.
    \item Regularly check and replace damaged gas pipes.
    \item If there is a gas leak, open windows and doors. Avoid using matchsticks, lighters and switches.
    \item Wear cotton clothes while cooking as cotton is less likely to catch fire, in case your clothes come in contract with fire.
    \item Turn off the main power supply in case of electric sparks.
    \item Regularly inspect and repair faulty electric points and frayed wires used in appliances.
\end{enumerate}
\subsection{\textbf{If paper, wood or plastic catch fire}}
If paper, wood or plastic catch fire, you can extinguish it using water. Use a bucket, a hose or any other available water source to douse the flames.
\subsection{\textbf{If a fire is caused by a short circuit}}
Never use water to extinguish an electrical fire. Water conducts electricity and can result in electrocution or the fire spreading. We may get a severe shock, which can even lead to death. We should switch OFF the main power supply.
\vspace{0.1in}
\par
We can use sand to put out electrical fires.
\subsection{\textbf{If a fire caused by fuels (petrol, diesel and kerosene)}}
Never use water to put out fires caused by fuels like petrol, diesel and kerosene. As oil is lighter than water, it keeps floating on the water and fire continues to burn. Instead, we should put sand or mud on the fire caused by fuels.
\vspace{0.4in}
\subsection{\textbf{If a person's clothes catch fire}}
If one's clothes catch fire, one should not try to run. Instead, one should STOP, DROP and ROLL. One should stop, drop to the ground, cover the face with with hands and roll over and over until the fire is out.
\vspace{0.1in}
\par
Once fire is out, immediately take the person to the hospital.
\vspace{2in}
\section{Internet Safety}
Internet safety is crucial or everyone.
\vspace{0.6in}
\subsection{Safety Tips}
\begin{enumerate}
    \item Always seek permisison from a parent or a guardian before using the Internet, especially for new websites or apps.
    \item Never share personal information online, such as your full name address, school name or passwords.
    \item Create strong passwords for online accounts, using a mix of letters, numbers and symbols.
    \item Don't click on links or pictures immediately; check with a grown-up to ensure they are safe to visit or view.
    \item Don't share your pictures without asking your parents first.
    \item Visit only approved websites, games and apps as designated by your parents or teachers.
    \item If something online makes you feel scared or confused, inform a trusted adult immediately.
    \item Balance online and offline activities. Spend time with family and friends.
    \item Be cautious when talking to unknown people online; avoid sharing personal information with them.
\end{enumerate}
\section{First Aid}
 Accidents can happen despite precautions. Knowing what to do is crucial. \textbf{First aid} is the assistance given to an injured person before professional help arrives. It aims to make the victim comfortable, ease pain and prevent the condition from getting worse. A small help and the right steps can sometimes save a life.
 \subsection{First Aid For Cuts And Wounds}
 First aid for cuts and wounds involves cleaning, treating and protecting the wound from infection and further damage.
 \begin{enumerate}
     \item Apply direct pressure on the affected area with a clean cloth or a gauze pad to stop bleeding.
     \item Once bleeding is controlled, gently clean the wound with mild soap and running water.
     \item Use a sterile gauze pad or a clean cloth to remove dirt without scrubbing the wound.
     \item Apply an antiseptic solution on the affected area to prevent infection.
     \item In case of heavy bleeding, use a tourniquet-a piece of cloth tightly wrapped around the world.
\end{enumerate}
\vspace{0.3in}
\par
I wonder if wounds should be left open or covered for quick healing!
\subsection{First Aid For Sprains}
Two bones are joined by a \textbf{ligament.} When a ligament is damaged, it causes a sprain. It leads to swelling and pain. For a sprain, follow these first aid steps:
\begin{enumerate}
     \item Ask the injured person to stop any activity and immediately rest the affected area.
     \item Avoid putting weight on the injured limb.
     \item Apply ice wrapped in a cloth for 15-20 minutes every 1-2 hours during the first 48 hours to reduce to swelling and pain.
     \item Wrap a crepe bandage around the injured area to control swelling.
     \item Elevate the injured area above the heart to further reduce swelling.
\end{enumerate}
\vspace{1.1in}
\subsection{First Aid For Fractures}
A \textbf{fracture} is a crack or a break in a bone. The first aid for fractures involves minimising the movement of the injured area to prevent further damage and reduce pain. The following first aid can be given for fractures:
\begin{enumerate}
    \item Keep the injured limb or area as still as possible to prevent further damage.
    \item If a fracture is suspected in an arm or a leg, make a splint using available materials (boards, rolled up newspapers or clothing). Secure the splint with bandages.
    \item If possible, elevate the injured slightly above the heart to reduce swelling.
    \item Support the injured area with cushions or pillows to keep it in a comfortable position.
    \item Apply an ice pack to reduce pain and swelling. Place a barrier between the ice pack and the skin.
\end{enumerate}
\vspace{0.7in}
\subsection{First Aid For Nosebleeds}
To stop a nosebleed, follow these steps:
\begin{enumerate}
    \item Lean slightly forward; do not tilt the head backward to prevent blood from flowing down the throat.
    \item Pinch the nostrils together just below the bony bridge of the nose using the thumb and the index finger.
    \item Keep pressing for 10-15 minutes without releasing. This helps to constrict blood vessels and stop bleeding. While holding the nostrils, breathe through the mouth.
    \item Use an ice pack wrapped in a cloth on the bridge of the nose to reduce bleeding.
\end{enumerate}
\subsection{First Aid For Minor Burns}
First aid for minor burns involves immediate actions to reduce pain and prevent infection. A minor burn may cause superficial redness of the skin and blisters.
\begin{enumerate}
    \item Hold burn under cool running water for about 10-15 minutes. This helps to reduce pain and prevent the burn from getting worse.
    \item Do not use ice as it can damage the skin and worsen the burn.
    \item Apply a thin layer of an antibiotic ointment to prevent infection.
    \item Cover the burn with a clean, sterile gauze to protect it from dirt.
\end{enumerate}
\vspace{0.1in}
\par
I wonder why ice should not be applied directly to a burn wound!
\vspace{0.4in}
\subsection{First Aid For Animal Bites}
Animal bites require promptly cleaning the wound, controlling bleeding and preventing infection. Here's what one should do:
\begin{enumerate}
    \item If the bite is bleeding, gently press a clean cloth on it.
    \item If possible, elevate the bitten area to reduce the blood flow.
    \item Clean the wound with a mild soap and warm water. Rinse it well and pat it dry with a clean cloth.
    \item Apply an antiseptic ointment to reduce infection risk.
    \item For bites of dogs, cats or monkeys (which can transmit rabies), seek immediate medical attention to get an anti-rabies injection.
    \item In the case of insect bites, like bee stings, remove the sting and wash the area thoroughly. Apply baking soda on bee bites and lemon on wasp bites.
\end{enumerate}
\vspace{1.1in}
\section{Now I Know}
\begin{itemize}
    \item It is a very important for all us to follow safety rules and stay safe.
    \item A fire may be many factors such as a fault in electric wires, leakage of cooking gas, bursting crackers or accidents while using fuels like petrol, kerosene and diesel.
    \item Internet safety is crucial for people all to protect their personal information and privacy.
    \item \textbf{First aid} refers to the immediate an initial assistance given to an injured person before professional medical help arrives.
\end{itemize}